% Active research-report draft. The untouched supplied template is preserved
% in official_template.tex; the previous six-page attempt is in previous_version/.
\documentclass[journal,10pt,twoside]{IEEEtran}

\usepackage[utf8]{inputenc}
\usepackage[T1]{fontenc}
\usepackage{graphicx}
\usepackage{amsmath,amssymb,mathtools}
\usepackage{booktabs}
\usepackage{tabularx}
\usepackage{enumitem}
\usepackage{xcolor}
\usepackage{balance}
\usepackage{microtype}
\usepackage{dblfloatfix}
\usepackage{url}
\usepackage[bookmarks=true,bookmarksopen=true]{hyperref}
\usepackage{cleveref}
\graphicspath{{images/}{../../docs/meetings/2026-10-19_update/figures/}}
\newcommand{\checkme}[1]{\textcolor{red}{[CHECK: #1]}}
\usepackage{natbib}\input{natbib-add}
\bibliographystyle{apalike}
\bibpunct{[}{]}{;}{a}{}{}

\setlist[itemize]{noitemsep,nolistsep,leftmargin=*}
\setlist[enumerate]{noitemsep,nolistsep,leftmargin=*}
\hypersetup{colorlinks=true,urlcolor=blue,linkcolor=blue,citecolor=blue}
\crefname{figure}{figure}{figures}
\Crefname{figure}{Figure}{Figures}
\crefname{table}{table}{tables}
\Crefname{table}{Table}{Tables}

\newcommand{\reporttitle}{Spatial Aware Peak Ranking in Mass Spectrometry Imaging}
\newcommand{\authorname}{Zayd \textsc{Suliman}}
\newcommand{\studentnumber}{2653934}
\newcommand{\supervisorname}{Hairong Bau and Richard Klein}
\newcommand{\academicyear}{2026}

\begin{document}

\title{%
  \vspace{4pt}
  \includegraphics[width=0.12\textwidth]{images/witslogonew.png}\\[6pt]
  \Large\textbf{\reporttitle}%
}

\author{%
  \authorname,~Student~No.~\studentnumber\\
  \small School of Computer Science and Applied Mathematics,
         Faculty of Science, University of the Witwatersrand\\
  \small Supervisors:~\supervisorname \qquad \academicyear
}

\maketitle

\begin{abstract}
Mass spectrometry imaging (MSI) measures thousands of mass-to-charge ($m/z$)
channels at every location of a tissue section, so analyses usually begin by
selecting a small set of informative peaks. Learned peak-selection methods such
as msiPL and S3PL rank $m/z$ bins from unsupervised models, and S3PL adds
spatial neighbourhood information through 3D convolution. This report asks
whether cheap local spatial context can give a lightweight variational
autoencoder (VAE) the same benefit, and which part of a learned pipeline
actually determines peak quality. We fitted VAE variants with and without
neighbourhood input, segmented their latent spaces with a Gaussian mixture
model (GMM), and ranked peaks either by Integrated Gradients (IG) through the
model or by direct correlation of each ion image with the GMM posterior maps.
All methods were compared at matched peak budgets on 24 sections from three
public collections (glioblastoma, colorectal adenocarcinoma and renal cell
carcinoma) under predeclared decision rules, against reproduced msiPL and S3PL
baselines and model-free rankings. Holding the segmentation fixed, posterior
correlation outperformed IG in 136 of 140 paired comparisons across three
training seeds, and this was statistically supported in all six
dataset--model comparisons (Holm-adjusted $p \le 0.047$). Neighbourhood input to
the encoder gave no consistent, spatially specific gain, whereas spatially
aware ranking statistics, namely segmentation correlation and Moran's~$I$,
were the strongest label-free rankings. A linear PCA segmentation performed
comparably to the VAE. Finally, S3PL's published glioblastoma score lay outside
the range of ten reproduced seeds; seeds altered S3PL mainly through weight
initialisation, which switched individual sections between distinct solutions.
The contribution is a controlled, statistically reported evaluation that bounds
where spatial information helps learned MSI peak ranking.
\end{abstract}

\begin{IEEEkeywords}
mass spectrometry imaging, peak selection, variational autoencoder, spatial context, Integrated Gradients, Moran's I, reproducibility
\end{IEEEkeywords}

\input{sections/introduction}
\input{sections/related_work}
\input{sections/proposal_outcomes}
\input{sections/methodology}
\input{sections/results}
\input{sections/discussion}
\input{sections/conclusion}

\section*{Acknowledgements}
The author thanks Hairong Bau and Richard Klein for their supervision and
guidance, and the University of the Witwatersrand for access to the MSCluster
computing facility.

\balance
\bibliography{references}

\appendices
\input{sections/appendix}
\end{document}
